\section{Introduction}
\label{sec:introduction}

Temporal learning is difficult for two related reasons. First, a prediction at
one time may depend on events far in the past, so the learner needs a state that
summarizes history without growing with the sequence length. Second, training
examples extracted from a single trajectory are generally dependent, so their
statistical behavior cannot be read directly from an i.i.d. analysis. Reservoir
computing addresses the first difficulty by driving a fixed dynamical system
with the input sequence and training only a comparatively simple readout
\citep{Jaeger2001,MaassNatschlagerMarkram2002}. When the state dynamics are
contractive, the effect of an arbitrary initialization vanishes, yielding a
well-defined causal representation of the input history.

Quantum reservoir computing replaces the classical recurrent state by a
quantum system and constructs classical features from measurements of the
evolved state \citep{FujiiNakajima2017,ChenNurdinYamamoto2020}. This offers a
structured way to generate nonlinear temporal features while keeping the
trainable component classical. It also makes the resource picture more
layered than in an abstract reservoir model. Input dimension determines how
many coordinates must be encoded, reservoir width and circuit depth determine
the cost of one state evolution, the observable family determines the feature
dimension, and finite-shot estimation determines how often the same temporal
window must be replayed. These quantities should be analyzed separately rather
than summarized by a single claim of scalability.

A second source of ambiguity is that several distinct questions are often
compressed into one performance claim. A contractive state update can guarantee
loss of initialization dependence without guaranteeing that different input
histories produce different features. A positive-definite kernel can
interpolate distinct feature vectors without establishing that the featurizer
is injective. Likewise, an efficient estimator for many local observables does
not by itself provide a learning guarantee under temporal dependence or a
model of device noise. A useful theory for quantum reservoir learning should
therefore connect these layers while preserving their individual assumptions.

We develop this perspective through \textsc{QuaRK}, a quantum-reservoir kernel
for temporal learning. Each input vector is first mapped to the width of an
\(n\)-qubit register by a finite-set Johnson--Lindenstrauss (JL) projection and
a bounded angle encoding. A collection of independently parameterized
reset-contractive quantum channels then processes a temporal window. Local
observable expectations from the resulting reservoir states are concatenated
into a classical feature vector, on which a regularized Mat\'ern kernel readout
is trained. The components have complementary roles: the projection controls
encoded width, contraction controls the influence of the remote past,
reservoir multiplexing and observable choice control feature dimension, and
RKHS regularization controls the readout class.

The analysis follows the same order as the learning pipeline. We first account
for the finite-set projection dimension, then use exact trace-norm contraction
to construct a unique causal state and quantify the error introduced by a
finite window. The geometry of the Mat\'ern RKHS transfers this state-space
error to predictions and squared loss. We next characterize the finite-sample
capacity of the readout conditional on distinct exact features. Finally, an
independent-block argument gives a one-sided risk bound for stationary
beta-mixing input--output data. We extend this result from one fixed design to
a predeclared finite family, so validation-based selection among finitely many
reservoir and readout configurations incurs an explicit logarithmic penalty.
A separate classical-shadow result quantifies the measurements needed to
estimate local Pauli features. Keeping model selection and measurement
separate from the exact-feature risk bound makes the scope of each result
explicit.

Throughout the paper, \emph{resource-aware} means that logical width, temporal
channel applications, state preparations, observable locality, classical
projection work, feature dimension, and kernel-training costs are reported as
separate resources. It is not used as a synonym for computational advantage or
practical efficiency.

\paragraph{Contributions.}
The main contributions are:
\begin{enumerate}
    \item We formulate an end-to-end temporal learner built from
    reset-contractive quantum reservoirs and a classical Mat\'ern kernel
    readout. Spatial multiplexing is represented by a tuple of independently
    evolved reservoir states followed by concatenated observable features,
    which keeps peak width and repeated execution costs distinct.

    \item We prove exact trace-norm contraction of the reset dynamics, derive a
    unique causal state and a geometrically decaying finite-window feature
    error, and propagate this approximation through the Mat\'ern RKHS to
    prediction and squared loss.

    \item We separate readout capacity from statistical generalization.
    Conditional on pairwise-distinct exact features, strict positive
    definiteness of the Mat\'ern kernel yields a finite-norm interpolant. For a
    fixed exact featurizer and stationary beta-mixing data, a corrected
    independent-block analysis yields a one-sided risk bound with an explicit
    effective block count and finite-window approximation term. A finite-family
    corollary makes the bound simultaneous over a predeclared collection of
    configurations and quantifies the additional model-selection cost.

    \item We state the local-Pauli classical-shadow requirement independently
    of the learning bound and derive how coordinate-wise measurement error
    affects evaluation of a fixed RKHS readout.
\end{enumerate}

The resulting contribution is an assumption-explicit foundation for studying
quantum-reservoir-based temporal features. The theory does not establish
quantum advantage or superiority over classical reservoirs; those are
comparative empirical questions that require matched protocols.

\paragraph{Organization.}
After positioning the work in \cref{sec:related-work}, we formalize temporal
windows, reservoir features, risks, and dependence in
\cref{sec:problem-setting}. \Cref{sec:method} then builds the complete learner
and accounts for its resources. The theoretical analysis is developed in
\cref{sec:theory}, with proofs in \cref{app:proofs}.  The empirical section
then specifies the controlled functional hierarchy, dependent-family
robustness study, high-dimensional width-control experiment, matched controls,
and measurement ablations.  The exact synthetic generators and task maps are
collected in \cref{app:dependent-families,app:functionals}, followed by the
discussion and conclusion.
